% ============================================================
%  cap05/5.10_limitaciones.tex
% ============================================================
\subsection{Limitaciones e Hilos Conductores Futuros}

A pesar de los abultados m\'eritos emp\'iricos registrados, \textit{Democra} no est\'a exenta de restricciones t\'ecnicas en su envoltura iterativa actual:
\begin{itemize}
    \item \textbf{Dependencia Conectiva Absoluta:} Al cimentarse \'integramente como un protocolo SaaS y SPA, la inoperatividad de los nodos de voluntariado en zonas rurales carentes de se\~nal 4G/5G sigue siendo una condici\'on letal. La integraci\'on futura de librer\'ias de \textit{Service Workers} (convirtiendo a Democra en una PWA con soporte \textit{Offline-First} mediante cach\'e IndexedDB) debe ser el pelda\~no subsiguiente.
    \item \textbf{Carga H\'ibrida de Documentos F\'isicos:} A\'un se requiere la digitalizaci\'on manual (escaneo de c\'odigos QR o fotograf\'ias) para los donativos s\'olidos (v\'iveres), lo que todav\'ia induce a micro-errores de error humano tipogr\'afico. El procesamiento automatizado de im\'agenes v\'ia redes neuronales (OCR \textit{Edge-Computing}) en los dispositivos m\'oviles aliviar\'ia esta carga asint\'otica.
\end{itemize}

El desarrollo actual, sin embargo, forja una cimentaci\'on sumamente resistente, s\'olida y l\'ogicamente as\'eptica sobre la cual se puede extender una m\'iriada de m\'odulos anal\'iticos sin comprometer el fr\'agil tejido de la soberan\'ia de datos inter-institucional.
